\chapter{Stochastic Calculus}\label{ch:iv1:stochastic-calculus}

\begin{ivtrains}
Brownian motion as a tool: Ito's lemma, martingales and optional stopping, reflection and hitting times, integrals of the path, Girsanov's change of measure and Feynman--Kac. Most results here can be reached in two ways, and the second is the check.
\end{ivtrains}

\begin{question}[$\star$ \iqroles{QR, QT, QD, Risk} \iqfirm{any} \iqlevel{junior}]\label{iq:iv1:stochastic-calculus:1}
$W$ is a standard Brownian motion. Give $\operatorname{Cov}(W_s, W_t)$. Find the variance of $W_1 + W_2$, the correlation of $W_1$ and $W_4$, and the probability that $W_1$ and $W_2$ are both positive.
\end{question}

\begin{question}[$\star$ \iqroles{QR, QT, QD, Risk} \iqfirm{any} \iqlevel{junior}]\label{iq:iv1:stochastic-calculus:2}
Which of these are martingales: $W_t$, $W_t^2$, $W_t^2 - t$, $tW_t$, $e^{W_t - t/2}$, $W_t^3$, $W_t^3 - 3tW_t$? Give one test that decides all of them.
\end{question}

\begin{question}[$\star$ \iqroles{QR, QT, QD, Risk} \iqfirm{bank, hedge fund} \iqlevel{junior}]\label{iq:iv1:stochastic-calculus:3}
A stock follows $dS = \mu S\,dt + \sigma S\,dW$ with $\mu = 10\%$ and $\sigma = 20\%$ a year. Use Ito's lemma on $\ln S$. Starting from 1, what are the mean and median after one year? What is the probability of ending above the mean?
\end{question}

\begin{question}[$\star$ \iqroles{QR, QT, QD, Risk} \iqfirm{any} \iqlevel{junior}]\label{iq:iv1:stochastic-calculus:4}
Compute $\mathbb{E}W_t^4$ and $\mathbb{E}e^{aW_t}$ from the normal distribution. What is $\mathbb{E}e^{2W_1}$? Find the correlation of $W_1^2$ with $W_4^2$.
\end{question}

\begin{question}[$\star$ \iqroles{QR, QT, QD, Risk} \iqfirm{any} \iqlevel{junior}]\label{iq:iv1:stochastic-calculus:5}
Show that $cW_{t/c^2}$ is a Brownian motion for any $c \ne 0$. Use this to relate $P(W_4 > 2)$ to a probability at time 1, and give its value. A return has a daily volatility of 1\%. What is the annual volatility, with 252 days and independent returns?
\end{question}

\begin{question}[$\star$ \iqroles{QR, QT, QD, Risk} \iqfirm{any} \iqlevel{junior}]\label{iq:iv1:stochastic-calculus:6}
Split $[0, 1]$ into 1000 equal steps. Find the mean and the standard deviation for the sum of squared increments of $W$, and what is the mean of the sum of the absolute increments? What happens to each as the steps shrink?
\end{question}

\begin{question}[$\star\star$ \iqroles{QR, QT, QD} \iqfirm{bank, hedge fund} \iqlevel{junior, experienced}]\label{iq:iv1:stochastic-calculus:7}
Find $\int_0^TW_t\,dW_t$ using Ito's lemma. Give its mean and variance. You approximate it by a sum over small steps, evaluating $W$ at the left end of each step, then at the right end, then at the midpoint. What does each sum converge to?
\end{question}

\begin{question}[$\star\star$ \iqroles{QR, QT, QD} \iqfirm{bank, hedge fund} \iqlevel{junior, experienced}]\label{iq:iv1:stochastic-calculus:8}
Let $I_T = \int_0^TW_t\,dt$. Find its distribution, using an integration by parts, and check the variance by a double integral. What is its correlation with $W_T$? Evaluate the variance at $T = 3$.
\end{question}

\begin{question}[$\star\star$ \iqroles{QR, QT, QD, Risk} \iqfirm{bank, hedge fund} \iqlevel{junior, experienced}]\label{iq:iv1:stochastic-calculus:9}
State the reflection principle, and derive the distribution of $M_T = \max_{s \le T}W_s$. What are $P(M_1 \ge 1)$ and $\mathbb{E}M_1$? How does your answer change if the barrier is only checked at daily closes, with 252 closes in the year?
\end{question}

\begin{question}[$\star\star$ \iqroles{QR, QT, QD, Risk} \iqfirm{any} \iqlevel{junior, experienced}]\label{iq:iv1:stochastic-calculus:10}
$T_a$ is the first time that $W$ reaches $a > 0$. Show that $T_a$ is finite with probability one, and find its median for $a = 1$. What is $P(T_1 > 100)$? Why is its mean infinite?
\end{question}

\begin{question}[$\star\star$ \iqroles{QR, QT, QD, Risk} \iqfirm{any} \iqlevel{junior}]\label{iq:iv1:stochastic-calculus:11}
$W$ starts at 0 and stops on first reaching $-1$ or $2$. Using two martingales, find the probability that it ends at 2 and the expected stopping time. What is the general formula for $-a$ and $b$? What happens as $a \to \infty$?
\end{question}

\begin{question}[$\star\star$ \iqroles{QR, QT, QD, Risk} \iqfirm{hedge fund, bank} \iqlevel{junior, experienced}]\label{iq:iv1:stochastic-calculus:12}
$X_t = \mu t + \sigma W_t$ with $\mu > 0$ and $X_0 = 0$. Find an exponential martingale, and use it to derive the probability that $X$ ever falls to $-a$. A strategy has an annual mean of 10\% and a volatility of 20\%, with no compounding. What is the probability that it ever loses 20\% from the start?
\end{question}

\begin{question}[$\star\star$ \iqroles{QR, QT, QD, Risk} \iqfirm{bank, hedge fund} \iqlevel{junior, experienced}]\label{iq:iv1:stochastic-calculus:13}
For $dS = \mu S\,dt + \sigma S\,dW$, apply Ito's lemma to $S^p$. Find $\mathbb{E}S_T^2$, and find the $p$ that makes $S^p$ a martingale, given $\mu = 10\%$ and $\sigma = 20\%$. Starting from 1, what is the probability of reaching 2 before $\frac12$?
\end{question}

\begin{question}[$\star\star$ \iqroles{QR, QT, QD, Risk} \iqfirm{bank, hedge fund} \iqlevel{junior, experienced}]\label{iq:iv1:stochastic-calculus:14}
Given $W_1 = x$, what are the conditional mean and variance of $W_s$, $0 < s < 1$? Given $W_1 = 0$, what is the probability that $W$ reaches level 1 during the interval? Derive the formula by reflection, for any $x$.
\end{question}

\begin{question}[$\star\star$ \iqroles{QR, QT, QD, Risk} \iqfirm{any} \iqlevel{junior, experienced}]\label{iq:iv1:stochastic-calculus:15}
Solve the Ornstein--Uhlenbeck equation $dX = -\kappa X\,dt + \sigma\,dW$. Give the mean and variance at $t$, the stationary variance and the half-life of a shock. Evaluate these at $\kappa = 2$, $\sigma = 1$, $t = 1$, $X_0 = 3$.
\end{question}

\begin{question}[$\star\star$ \iqroles{QR, QT, QD, Risk} \iqfirm{bank, hedge fund} \iqlevel{junior}]\label{iq:iv1:stochastic-calculus:16}
$W^1$ and $W^2$ have correlation $\rho$. Write one in terms of the other and an independent motion. Use Ito's lemma to find $\mathbb{E}W^1_tW^2_t$. At $\rho = 0.5$ and $t = 1$, find $\operatorname{Var}(W^1_t + W^2_t)$ and $\mathbb{E}[(W^1_t)^2(W^2_t)^2]$.
\end{question}

\begin{question}[$\star\star$ \iqroles{QR, QT, QD} \iqfirm{bank, hedge fund} \iqlevel{junior, experienced}]\label{iq:iv1:stochastic-calculus:17}
Show directly, by completing the square, that $\mathbb{E}[f(W_T + \mu T)] = \mathbb{E}[f(W_T)e^{\mu W_T - \mu^2T/2}]$. Use the identity to estimate $P(W_1 > 4)$ by simulation, with a tilted drift. By how much does this reduce the standard error, for one million paths?
\end{question}

\begin{question}[$\star\star$ \iqroles{QR, QT, QD, Risk} \iqfirm{bank, hedge fund} \iqlevel{junior}]\label{iq:iv1:stochastic-calculus:18}
A stock has $\mu = 10\%$, $\sigma = 20\%$, $S_0 = 100$, and the rate is zero. A one-year digital pays 1 if $S_1 > 100$. What are the real-world probability, the risk-neutral probability and the price? What is the market price of risk, and how does it relate the two measures? Does the price depend on $\mu$?
\end{question}

\begin{question}[$\star\star$ \iqroles{QR, QT, QD} \iqfirm{bank, hedge fund} \iqlevel{junior, experienced}]\label{iq:iv1:stochastic-calculus:19}
State the Feynman--Kac formula for $u_t + \mu u_x + \frac12\sigma^2u_{xx} - ru = 0$, with terminal condition $u(T, x) = g(x)$. Solve the equation with $\mu = r = 0$, $\sigma = 1$, first for $g = x^2$, then for $g = e^x$. Check each by substituting into the equation. How does a rate $r$ change them?
\end{question}

\begin{question}[$\star\star$ \iqroles{QR, QT, QD, Risk} \iqfirm{hedge fund} \iqlevel{junior, experienced}]\label{iq:iv1:stochastic-calculus:20}
Let $X_t = e^{W_t - t/2}$. What equation does it solve? What are $\mathbb{E}X_T$, the median and $P(X_T < 1)$ at $T = 1$ and $T = 100$? What is the long-run behaviour of a single path? How do you reconcile this with the mean?
\end{question}

\begin{question}[$\star\star\star$ \iqroles{QR, QT, QD} \iqfirm{bank, hedge fund} \iqlevel{experienced}]\label{iq:iv1:stochastic-calculus:21}
What are the conditional mean and variance of $\int_0^TW_t\,dt$ given $W_T$? Evaluate them for $T = 1$. What share of the variance of the integral is explained by the endpoint?
\end{question}

\begin{question}[$\star\star\star$ \iqroles{QR, QT} \iqfirm{hedge fund} \iqlevel{junior, experienced}]\label{iq:iv1:stochastic-calculus:22}
Let $A$ be the fraction of $[0, 1]$ that $W$ spends above 0. State the distribution of $A$. Give the mean, the variance and $P(A \le 0.1)$. A trader's cumulative profit behaves like $W$. What does this say about how often a skill-free trader is ahead for at least 90\% of the year?
\end{question}

\begin{question}[$\star\star\star$ \iqroles{QR, QT} \iqfirm{hedge fund} \iqlevel{junior, experienced}]\label{iq:iv1:stochastic-calculus:23}
Let $M_T$ be the running maximum of $W$. From the reflection principle, find $P(M_T \ge a, W_T \le x)$ for $x \le a$. Show that $M_T - W_T$ has the same law as $|W_T|$. A driftless strategy has a volatility of 20\% with no compounding. What is the expected distance below its high after a year, and what is the chance that this exceeds 20\%?
\end{question}

\begin{question}[$\star\star\star$ \iqroles{QR, QT, QD, Risk} \iqfirm{hedge fund, bank} \iqlevel{experienced}]\label{iq:iv1:stochastic-calculus:24}
$X_t = \mu t + \sigma W_t$ with $\mu > 0$, and $\tau$ is the first time that $X$ reaches $a > 0$. Use martingales to derive $\mathbb{E}\tau$ and $\operatorname{Var}\tau$. Evaluate them at $a = 1$, $\mu = 0.1$ and $\sigma = 0.2$. What is $P(\tau > 20)$, and what distribution is this?
\end{question}

\begin{question}[$\star\star\star$ \iqroles{QR, QT} \iqfirm{hedge fund} \iqlevel{experienced}]\label{iq:iv1:stochastic-calculus:25}
Use an exponential martingale and optional stopping to show that $\mathbb{E}e^{-\lambda T_a} = e^{-a\sqrt{2\lambda}}$, where $T_a$ is the first time $W$ reaches $a > 0$. Why must you choose the positive root? What do you learn by letting $\lambda \to 0$, and from the derivative there? Evaluate at $a = 1$ and $\lambda = \frac12$.
\end{question}

\begin{question}[$\star\star\star$ \iqroles{QR, QT, QD, Risk} \iqfirm{hedge fund, bank} \iqlevel{experienced}]\label{iq:iv1:stochastic-calculus:26}
An asset follows geometric Brownian motion with volatility 20\%. Starting from 1, find the expected time to reach 2 if $\mu = 10\%$. If $\mu = 1\%$, what is the probability that it ever reaches 2? What happens at $\mu = 2\%$? Check the second result with a martingale of the form $S^p$.
\end{question}

\begin{question}[$\star\star\star$ \iqroles{QR, QT, QD} \iqfirm{hedge fund, bank} \iqlevel{experienced}]\label{iq:iv1:stochastic-calculus:27}
Ito's lemma does not apply to $|x|$, which has a kink at 0. Using Tanaka's formula, write $|W_t|$ as a stochastic integral plus a local time $L_t$. Find $\mathbb{E}L_1$. What does $L$ measure, and how could you estimate it from a simulated path? What other process has the same law as $L_t$?
\end{question}

\begin{question}[$\star\star\star$ \iqroles{QR, QT} \iqfirm{hedge fund} \iqlevel{experienced}]\label{iq:iv1:stochastic-calculus:28}
Find the mean and variance of $Q = \int_0^1W_t^2\,dt$ by double integrals. Then check both using the expansion $W_t = \sum_k\sqrt{2}\,Z_k\sin((k - \frac12)\pi t)/((k - \frac12)\pi)$, with independent standard normal $Z_k$, using $\sum_k((k - \frac12)\pi)^{-2} = \frac12$. What sum must you use for the variance?
\end{question}

\begin{question}[$\star\star\star\star$ \iqroles{QR, QT} \iqfirm{hedge fund} \iqlevel{experienced}]\label{iq:iv1:stochastic-calculus:29}
Let $Q = \int_0^1W_t^2\,dt$. Derive $\mathbb{E}e^{-\theta^2Q/2} = 1/\sqrt{\cosh\theta}$ using the Feynman--Kac formula, with the guess $u(t, x) = \exp(-\frac12A(t)x^2 + B(t))$. Solve for $A$ and $B$. Check the small-$\theta$ limit against the mean of $Q$, which is $\frac12$. Evaluate it at $\theta = 1$.
\end{question}

\begin{question}[$\star\star\star\star$ \iqroles{QR, QT} \iqfirm{hedge fund} \iqlevel{experienced}]\label{iq:iv1:stochastic-calculus:30}
$\tau$ is the first time that $W$ leaves $(-a, a)$. Using $\cosh(\theta W_t)e^{-\theta^2t/2}$, find $\mathbb{E}e^{-\lambda\tau}$. Expand it to obtain $\mathbb{E}\tau$ and $\operatorname{Var}\tau$. Check these using the polynomial martingale $W_t^4 - 6tW_t^2 + 3t^2$. Then use the poles of the transform to give the rate at which $P(\tau > t)$ decays.
\end{question}
